% !TeX program = pdflatex
\documentclass[aspectratio=169,xcolor={dvipsnames,table}]{beamer}

% Preámbulo modular para presentaciones Beamer (Modelación Estadística)
% Diseño y paleta institucional del Tecnológico de Monterrey

\documentclass[aspectratio=169,xcolor={dvipsnames,table}]{beamer}

\usepackage[utf8]{inputenc}
\usepackage[spanish,mexico]{babel}
\usepackage{amsmath,amssymb,amsfonts,mathtools}
\usepackage{graphicx}
\usepackage{listings}
\usepackage{booktabs}
\usepackage{tikz}
\usetikzlibrary{matrix,arrows,decorations.pathmorphing,shapes.geometric,positioning}

% Paleta Institucional Tecnológico de Monterrey
\definecolor{TecAzul}{HTML}{0039A6}       % Azul TEC: color primario institucional
\definecolor{TecAzulOscuro}{HTML}{1A2E51} % Azul marino oscuro
\definecolor{TecRojo}{HTML}{EC2661}       % Rojo de acento (paleta miTec)
\definecolor{TecGris}{HTML}{646464}       % Gris neutro para comentarios y subtítulos
\definecolor{TecFondoBloque}{HTML}{F4F6F9}% Fondo suave para bloques

% Configuración de Tema Beamer
\usetheme{Boadilla}
\usecolortheme[named=TecAzulOscuro]{structure}
\setbeamercolor{alerted text}{fg=TecRojo}
\setbeamercolor{palette primary}{bg=TecAzulOscuro,fg=white}
\setbeamercolor{palette secondary}{bg=TecAzul,fg=white}
\setbeamercolor{palette tertiary}{bg=TecAzulOscuro!80!black,fg=white}
\setbeamercolor{title}{fg=TecAzulOscuro,bg=TecFondoBloque}
\setbeamercolor{frametitle}{fg=TecAzulOscuro,bg=TecFondoBloque}

% Bloques personalizados elegantes
\setbeamercolor{block title}{bg=TecAzulOscuro,fg=white}
\setbeamercolor{block body}{bg=TecFondoBloque,fg=black}
\setbeamercolor{block title alerted}{bg=TecRojo,fg=white}
\setbeamercolor{block body alerted}{bg=TecRojo!8,fg=black}
\setbeamercolor{block title example}{bg=TecAzul,fg=white}
\setbeamercolor{block body example}{bg=TecAzul!8,fg=black}

% Redondear bordes y añadir sombra sutil
\setbeamertemplate{blocks}[rounded][shadow=true]
\setbeamertemplate{navigation symbols}{}

% Pie de página limpio con número de diapositiva
\setbeamertemplate{footline}{
  \leavevmode%
  \hbox{%
  \begin{beamercolorbox}[wd=.7\paperwidth,ht=2.25ex,dp=1ex,left]{author in head/foot}%
    \hspace*{2ex}\insertshortauthor~~(\insertshortinstitute) \hfill \insertshorttitle
  \end{beamercolorbox}%
  \begin{beamercolorbox}[wd=.3\paperwidth,ht=2.25ex,dp=1ex,right]{date in head/foot}%
    \insertshortdate{}\hspace*{2em}
    \insertframenumber{} / \inserttotalframenumber\hspace*{2ex}
  \end{beamercolorbox}}%
  \vskip0pt%
}

% Configuración de Listings de Python para visualización pedagógica de código
\lstset{
    language=Python,
    basicstyle=\ttfamily\footnotesize,
    keywordstyle=\color{TecAzulOscuro}\bfseries,
    stringstyle=\color{TecRojo},
    commentstyle=\color{TecGris}\itshape,
    showstringspaces=false,
    breaklines=true,
    frame=lines,
    rulecolor=\color{TecAzulOscuro!30},
    backgroundcolor=\color{TecFondoBloque},
    aboveskip=0.5em,
    belowskip=0.5em,
    literate={á}{{\'a}}1 {é}{{\'e}}1 {í}{{\'i}}1 {ó}{{\'o}}1 {ú}{{\'u}}1
             {Á}{{\'A}}1 {É}{{\'E}}1 {Í}{{\'I}}1 {Ó}{{\'O}}1 {Ú}{{\'U}}1
             {ñ}{{\~n}}1 {Ñ}{{\~N}}1 {¿}{{?`}}1 {¡}{{!`}}1
}

% Metadatos institucionales por defecto
\title[Modelación Estadística]{Modelación Estadística para Ciencia de Datos e Ingeniería}
\author[J. Castillo Colmenares]{Juliho David Castillo Colmenares}
\institute[Tec de Monterrey]{Tecnológico de Monterrey \\ Escuela de Ingeniería y Ciencias}
\date{\today}

% Comandos y macros estadísticas compartidas para las presentaciones Beamer (Metropolis)

% Conjuntos numéricos básicos
\newcommand{\R}{\mathbb{R}}
\newcommand{\N}{\mathbb{N}}
\newcommand{\Q}{\mathbb{Q}}
\newcommand{\Z}{\mathbb{Z}}
\newcommand{\set}[1]{\left\{ #1 \right\}}
\newcommand{\abs}[1]{\left|#1\right|}
\newcommand{\norm}[1]{\left\| #1 \right\|}

% Comandos estadísticos y probabilísticos del curso
\newcommand{\Var}{\operatorname{Var}}
\newcommand{\cov}{\operatorname{Cov}}
\newcommand{\corr}{\rho}
\newcommand{\comb}[2]{\begin{pmatrix} #1 \\ #2 \end{pmatrix}}
\newcommand{\s}{\sigma}
\newcommand{\particion}{\mathfrak{P}}
\newcommand{\card}[1]{\#\left( #1 \right)}
\newcommand{\seq}[3]{\set{#1}_{#2}^{#3}}
\newcommand{\E}{\mathbb{E}}
\newcommand{\Prob}{\mathbb{P}}

% Entornos pedagógicos y cajas en español académico natural
\newenvironment{discusion}{%
  \begin{alertblock}{Pregunta de Discusión}%
}{%
  \end{alertblock}%
}

\newenvironment{reflexion}{%
  \begin{alertblock}{Reflexión para el Aula}%
}{%
  \end{alertblock}%
}

\newenvironment{paradoja}{%
  \begin{alertblock}{Paradoja de Exploración}%
}{%
  \end{alertblock}%
}

\newenvironment{motivacion}{%
  \begin{exampleblock}{Motivación: Ciencia de Datos en la Práctica}%
}{%
  \end{exampleblock}%
}

\newenvironment{retoclase}{%
  \begin{block}{Reto Rápido en Equipo}%
}{%
  \end{block}%
}


\title{Funciones de Probabilidad Discretas}
\subtitle{Sección 03.02 --- Función de Distribución Acumulada (CDF Discreta)}
\author[J. Castillo Colmenares]{Juliho Castillo Colmenares}
\institute[Tec de Monterrey]{Tecnológico de Monterrey}
\date{}

\begin{document}

% SLIDE 01: Portada
\begin{frame}[plain]
	\titlepage
\end{frame}

% SLIDE 02: Hoja de Ruta
\begin{frame}{Hoja de Ruta --- Capítulo 03: Variables Aleatorias Discretas}
	\begin{block}{Estructura Curricular de la Unidad 2}
		\small
		\begin{enumerate}
			\itemsep=0.04cm
			\item \textcolor{gray}{Sección 03.01: Funciones de Probabilidad Discretas (PMF y Soporte)}
			\item \textbf{\textcolor{TecRojo}{Sección 03.02: Función de Distribución Acumulada (CDF Discreta)}} $\leftarrow$ \emph{Foco actual}
			\item \textcolor{gray}{Sección 03.03: Esperanza Matemática, Varianza y Momentos}
			\item \textcolor{gray}{Sección 03.04: Distribuciones Especiales I (Binomial, Geométrica e Hipergeométrica)}
			\item \textcolor{gray}{Sección 03.05: Distribuciones Especiales II (Poisson, Binomial Negativa y Multinomial)}
		\end{enumerate}
	\end{block}
	\vspace{-0.15cm}
	\begin{alertblock}{Objetivo Didáctico}
		\footnotesize
		Dominar la fundamentación analítica, geométrica y computacional de la Función de Distribución Acumulada $F(x)$ para variables discretas, caracterizando sus mesetas de acumulación, saltos de discontinuidad y operadores de inversión en el cálculo de probabilidades de intervalos y cuantiles.
	\end{alertblock}
\end{frame}

% SLIDE 03: Motivación e Intuición
\begin{frame}{Motivación --- ¿Por qué Acumular Probabilidad?}
	\begin{columns}[T]
		\begin{column}{0.48\textwidth}
			\begin{block}{Limitación de Consultas Puntuales}
				\small
				\begin{itemize}
					\itemsep=0.06cm
					\item La PMF $f(x) = P(X=x)$ responde únicamente a eventos de \emph{igualdad exacta}.
					\item En aplicaciones de ingeniería, gestión de riesgo y confiabilidad, las preguntas críticas son de \textbf{umbral} o de \textbf{intervalo}:
					\begin{itemize}
						\item \scriptsize ¿Probabilidad de que el servidor falle en $t \le 5$ ms?
						\item \scriptsize ¿Probabilidad de tener a lo sumo $3$ paquetes perdidos?
					\end{itemize}
				\end{itemize}
			\end{block}
		\end{column}
		\begin{column}{0.48\textwidth}
			\begin{block}{El Operador Acumulativo}
				\small
				\begin{itemize}
					\itemsep=0.06cm
					\item En lugar de sumar masas puntuales repetidamente, definimos una función global $F(x)$ que almacena la \emph{historia probabilística total} hasta el punto $x$.
					\item Esta función transforma el cálculo de sumatorias en simples evaluaciones y restas algebraicas en los extremos de interés.
				\end{itemize}
			\end{block}
		\end{column}
	\end{columns}
\end{frame}

% SLIDE 04: Consultas de Umbral vs. Puntuales
\begin{frame}{Consultas de Umbral vs. Consultas Puntuales}
	\begin{table}
		\centering
		\scriptsize
		\begin{tabular}{p{3.1cm}|p{4.5cm}|p{4.7cm}}
			\hline
			\textbf{Característica} & \textbf{Función de Masa (PMF) $f(x)$} & \textbf{Función Acumulada (CDF) $F(x)$} \\
			\hline \hline
			\textbf{Definición formal} & $f(x) = P(X = x)$ & $F(x) = P(X \le x) = \sum_{u \le x} f(u)$ \\
			\hline
			\textbf{Pregunta canónica} & ¿Exactamente cuántos? & ¿A lo sumo cuántos? \\
			\hline
			\textbf{Dominio e imagen} & $\mathbb{R} \to [0, 1]$, nula fuera de $\mathcal{S}_X$ & $\mathbb{R} \to [0, 1]$, definida $\forall x \in \mathbb{R}$ \\
			\hline
			\textbf{Geometría gráfica} & Diagrama de barras o bastones & Función escalonada (mesetas y saltos) \\
			\hline
			\textbf{Condición global} & $\sum_{x \in \mathcal{S}_X} f(x) = 1$ & $\lim_{x \to -\infty} F(x) = 0$, $\lim_{x \to \infty} F(x) = 1$ \\
			\hline
		\end{tabular}
	\end{table}
	\vspace{-0.2cm}
	\begin{alertblock}{Síntesis Conceptual}
		\footnotesize
		La CDF encapsula toda la información probabilística del sistema: conocer $F(x)$ para todo $x \in \mathbb{R}$ es estrictamente equivalente a conocer la distribución completa de $X$.
	\end{alertblock}
\end{frame}

% SLIDE 05: Definición Formal de la CDF
\begin{frame}{Definición Formal de la Función de Distribución Acumulada}
	\begin{block}{Definición Analítica (Ecuación 2.6.2 del Libro Maestro)}
		\scriptsize
		La \textbf{Función de Distribución Acumulada} (CDF) de una variable aleatoria discreta $X$ se define para todo $x \in \mathbb{R}$ como:
		\vspace{-0.15cm}
		\begin{align*}
			F(x) = P(X \le x) = \sum_{u \le x} f(u)
		\end{align*}
		\vspace{-0.25cm}
		donde la sumatoria se extiende sobre todos los elementos del soporte $u \in \mathcal{S}_X$ con $u \le x$.
	\end{block}
	\vspace{-0.4cm}
	\begin{columns}[T]
		\begin{column}{0.48\textwidth}
			\begin{exampleblock}{Soporte Finito Ordenado}
				\scriptsize
				Si $x_1 < x_2 < \dots < x_n$, entonces:
				\vspace{-0.15cm}
				\begin{align*}
					F(x) = \begin{cases}
						0 & x < x_1 \\
						f(x_1) & x_1 \le x < x_2 \\
						f(x_1) + f(x_2) & x_2 \le x < x_3 \\
						\vdots & \vdots \\
						1 & x \ge x_n
					\end{cases}
				\end{align*}
			\end{exampleblock}
		\end{column}
		\begin{column}{0.48\textwidth}
			\begin{alertblock}{Interpretación Geométrica}
				\scriptsize
				\begin{itemize}
					\itemsep=0.02cm
					\item $F(x)$ avanza por \textbf{mesetas horizontales} sin masa.
					\item En cada punto $x_k \in \mathcal{S}_X$, la función da un \textbf{salto vertical} de magnitud $f(x_k)$.
				\end{itemize}
			\end{alertblock}
		\end{column}
	\end{columns}
\end{frame}

% SLIDE 06: Propiedades Fundamentales I (Monotonía y Continuidad)
\begin{frame}{Propiedades Fundamentales --- Monotonía y Continuidad por la Derecha}
	\begin{block}{Propiedad 1: Monotonía No Decreciente}
		\footnotesize
		Para cualesquiera dos números reales $a, b \in \mathbb{R}$ tales que $a \le b$, se cumple rigurosamente que:
		\begin{align*}
			a \le b \implies \set{\omega : X(\omega) \le a} \subseteq \set{\omega : X(\omega) \le b} \implies F(a) \le F(b)
		\end{align*}
		La acumulación de probabilidades jamás puede decrecer, ya que las masas puntuales son no negativas ($f(x) \ge 0$).
	\end{block}
	\vspace{-0.15cm}
	\begin{block}{Propiedad 2: Continuidad por la Derecha (cadlág)}
		\footnotesize
		En todo punto $x \in \mathbb{R}$, el límite por la derecha coincide exactamente con el valor evaluado de la función:
		\begin{align*}
			F(x^+) = \lim_{u \to x^+} F(u) = F(x)
		\end{align*}
		En los puntos de salto $x_k \in \mathcal{S}_X$, el límite por la izquierda es estrictamente menor: $F(x_k^-) = F(x_k) - f(x_k)$. Por convención matemática universal, el punto superior del escalón se incluye (cierre por la derecha).
	\end{block}
\end{frame}

% SLIDE 07: Propiedades Fundamentales II (Límites Asintóticos)
\begin{frame}{Propiedades Fundamentales --- Límites Asintóticos y Normalización}
	\begin{block}{Propiedad 3: Condiciones de Frontera Asintóticas}
		\small
		Toda función de distribución acumulada legítima debe satisfacer las dos condiciones de contorno en el infinito:
		\begin{align*}
			\lim_{x \to -\infty} F(x) = 0 \quad \text{y} \quad \lim_{x \to \infty} F(x) = 1
		\end{align*}
	\end{block}
	\vspace{-0.15cm}
	\begin{columns}[T]
		\begin{column}{0.48\textwidth}
			\begin{alertblock}{Límite Inferior ($\to -\infty$)}
				\footnotesize
				A medida que retrocedemos hacia la izquierda de la recta real, dejamos atrás todos los puntos del soporte $\mathcal{S}_X$. La masa acumulada tiende a cero (evento imposible).
			\end{alertblock}
		\end{column}
		\begin{column}{0.48\textwidth}
			\begin{exampleblock}{Límite Superior ($\to \infty$)}
				\footnotesize
				Al avanzar infinitamente hacia la derecha, englobamos todo el espacio muestral $\Omega$. La suma total de los saltos converge exactamente a la normalización unitaria $1$.
			\end{exampleblock}
		\end{column}
	\end{columns}
\end{frame}

% SLIDE 08: Teorema de Inversión (Recuperación de la PMF)
\begin{frame}{Teorema de Inversión --- Recuperación de la PMF desde la CDF}
	\begin{block}{Teorema de Diferencia Finita (Ecuación 2.6.4 del Libro Maestro)}
		\small
		La función de masa de probabilidad $f(x)$ puede reconstruirse de manera única a partir de la función de distribución acumulada $F(x)$ mediante la magnitud exacto-algebraica del salto en el punto $x$:
		\begin{align*}
			f(x) = F(x) - \lim_{u \to x^-} F(u) \equiv F(x) - F(x^-)
		\end{align*}
	\end{block}
	\vspace{-0.15cm}
	\begin{columns}[T]
		\begin{column}{0.48\textwidth}
			\begin{exampleblock}{Soporte Discreto Ordenado}
				\scriptsize
				Si $\mathcal{S}_X = \set{x_1, x_2, \dots, x_n}$ con $x_1 < x_2 < \dots$:
				\begin{align*}
					f(x_1) &= F(x_1) - 0 = F(x_1), \\
					f(x_k) &= F(x_k) - F(x_{k-1}) \quad (k > 1).
				\end{align*}
			\end{exampleblock}
		\end{column}
		\begin{column}{0.48\textwidth}
			\begin{alertblock}{Puntos Fuera del Soporte}
				\scriptsize
				Si $x \notin \mathcal{S}_X$, la función $F(u)$ es localmente constante alrededor de $x$ (está en el interior de una meseta horizontal). Por tanto:
				\begin{align*}
					F(x) = F(x^-) \implies f(x) = 0.
				\end{align*}
			\end{alertblock}
		\end{column}
	\end{columns}
\end{frame}

% SLIDE 09: Álgebra de Probabilidades en Intervalos
\begin{frame}{Álgebra de Probabilidades en Intervalos}
	\begin{block}{Cálculo en Intervalos Semiabiertos por la Izquierda}
		\footnotesize
		Para cualquier par de números reales $a, b \in \mathbb{R}$ con $a \le b$, la probabilidad de que la variable caiga en $(a, b]$ es la diferencia exacta de sus acumulaciones:
		\begin{align*}
			P(a < X \le b) = P(X \le b) - P(X \le a) = F(b) - F(a)
		\end{align*}
	\end{block}
	\vspace{-0.15cm}
	\begin{block}{Tipología Completa de Intervalos para Variables Discretas}
		\scriptsize
		En variables discretas, la inclusión o exclusión de los extremos de frontera altera el resultado según la presencia de masa en $a$ y $b$:
		\begin{table}
			\centering
			\begin{tabular}{l|l|l}
				\hline
				\textbf{Tipo de Intervalo} & \textbf{Expresión de Probabilidad} & \textbf{Descomposición en CDF} \\
				\hline \hline
				Semiabierto $(a, b]$ & $P(a < X \le b)$ & $F(b) - F(a)$ \\
				\hline
				Cerrado $[a, b]$ & $P(a \le X \le b) = P(X=a) + P(a < X \le b)$ & $F(b) - F(a^-)$ \\
				\hline
				Abierto $(a, b)$ & $P(a < X < b) = P(a < X \le b) - P(X=b)$ & $F(b^-) - F(a)$ \\
				\hline
				Semiabierto $[a, b)$ & $P(a \le X < b) = P(a \le X \le b) - P(X=b)$ & $F(b^-) - F(a^-)$ \\
				\hline
			\end{tabular}
		\end{table}
	\end{block}
\end{frame}

% SLIDE 10: Cuantiles Discretos y la Mediana
\begin{frame}{Cuantiles Discretos e Inversión Generalizada}
	\begin{block}{Definición de Cuantil de Orden $\alpha$}
		\small
		Dado un nivel de probabilidad $\alpha \in (0, 1)$, el \textbf{cuantil discreto} $q_\alpha$ se define formalmente como el ínfimo de los números reales donde el acumulado alcanza o sobrepasa el umbral $\alpha$:
		\begin{align*}
			q_\alpha = \inf \set{x \in \mathbb{R} : F(x) \ge \alpha}
		\end{align*}
	\end{block}
	\vspace{-0.15cm}
	\begin{columns}[T]
		\begin{column}{0.48\textwidth}
			\begin{alertblock}{La Mediana ($q_{0.50}$)}
				\footnotesize
				Es el valor $x_{\text{med}}$ del soporte que divide a la distribución en dos mitades equilibradas:
				\begin{align*}
					P(X \le x_{\text{med}}) \ge 0.50 \text{ y } P(X \ge x_{\text{med}}) \ge 0.50
				\end{align*}
			\end{alertblock}
		\end{column}
		\begin{column}{0.48\textwidth}
			\begin{exampleblock}{Inversa Generalizada}
				\footnotesize
				Como las funciones escalonadas no son biyectivas (tienen tramos planos), la función ínfimo opera como una pseudo-inversa única $F^{-1}(\alpha)$ que mapea probabilidades a valores del soporte.
			\end{exampleblock}
		\end{column}
	\end{columns}
\end{frame}

% SLIDE 11: Lab Python --- Parte 1: Construcción de CDF
\begin{frame}[fragile]{Lab Computacional --- Parte 1: Construcción de la CDF}
	\begin{block}{Acumulación Programática de Probabilidades (\texttt{compute\_step\_cdf})}
		\scriptsize
		El siguiente bloque en Python ordena el soporte vectorial y aplica la suma acumulativa (\texttt{np.cumsum}) para generar las alturas exactas de las mesetas escalonadas de $F(x)$.
	\end{block}
	\vspace{-0.1cm}
	\lstinputlisting[language=Python, firstline=9, lastline=21, basicstyle=\fontsize{6.3pt}{7.5pt}\ttfamily]{../../code/03_variables_aleatorias_discretas/03.02_discrete_cdf.py}
\end{frame}

% SLIDE 12: Lab Python --- Parte 2: Inversión de Saltos PMF
\begin{frame}[fragile]{Lab Computacional --- Parte 2: Inversión de Saltos y PMF}
	\begin{block}{Algoritmo de Diferencia Finita (\texttt{extract\_pmf\_from\_cdf})}
		\scriptsize
		Mediante el operador de diferencia hacia atrás (\texttt{np.diff} con \texttt{prepend=0.0}), el script reconstruye exactamente la masa puntual de cada salto y verifica el error de normalización.
	\end{block}
	\vspace{-0.1cm}
	\lstinputlisting[language=Python, firstline=23, lastline=35, basicstyle=\fontsize{6.3pt}{7.5pt}\ttfamily]{../../code/03_variables_aleatorias_discretas/03.02_discrete_cdf.py}
\end{frame}

% SLIDE 13: Lab Python --- Parte 3: Probabilidad en Intervalos y Cuantiles
\begin{frame}[fragile]{Lab Computacional --- Parte 3: Consultas de Intervalo y Cuantiles}
	\begin{block}{Evaluación de Subconjuntos y Percentiles (\texttt{evaluate\_interval\_and\_quantiles})}
		\scriptsize
		El código utiliza búsqueda binaria eficiente (\texttt{np.searchsorted}) para restar acumulaciones en intervalos y determina la mediana y el percentil 85 inspeccionando umbrales lógicos.
	\end{block}
	\vspace{-0.1cm}
	\lstinputlisting[language=Python, firstline=37, lastline=52, basicstyle=\fontsize{6.3pt}{7.5pt}\ttfamily]{../../code/03_variables_aleatorias_discretas/03.02_discrete_cdf.py}
\end{frame}

% SLIDE 14: Lab Python --- Parte 4: Simulación y Convergencia ECDF
\begin{frame}[fragile]{Lab Computacional --- Parte 4: Simulación Monte Carlo y ECDF}
	\begin{block}{Convergencia de la Distribución Empírica (\texttt{simulate\_empirical\_ecdf})}
		\scriptsize
		Se simulan $N=100,000$ lanzamientos de dos dados para construir la ECDF $F_n(x)$ y comprobar por el Teorema de Glivenko-Cantelli una desviación máxima inferior a $0.005$.
	\end{block}
	\vspace{-0.2cm}
	\lstinputlisting[language=Python, firstline=54, lastline=73, basicstyle=\fontsize{5.4pt}{6.5pt}\ttfamily]{../../code/03_variables_aleatorias_discretas/03.02_discrete_cdf.py}
\end{frame}

% SLIDE 15: Problemas --- Nivel Fundamental
\begin{frame}{Problemas del Libro Maestro --- Nivel Fundamental}
	\begin{block}{Ejercicios Directos de Conceptualización (Taxonomía 3-3-2-2)}
		\small
		\begin{itemize}
			\itemsep=0.08cm
			\item \textbf{Problema 3.2.1 (Dado justo y saltos uniformes):} Construcción explícita en siete intervalos disjuntos de $F(x)$ para un dado equilibrado y evaluación numérica exacta en puntos ordinarios ($2.8$), saltos de frontera ($4$) y colas asintóticas ($-1, 10$).
			\item \textbf{Problema 3.2.2 (Reconstrucción desde tabla escalonada):} Cálculo de probabilidades en intervalos semiabiertos $P(0 < Y \le 5) = F(5) - F(0) = 0.70$ y acotados cerrados $P(2 \le Y \le 8) = 0.85$, deduciendo la tabla PMF por diferencias finitas en cada salto.
			\item \textbf{Problema 3.2.3 (Demostración de monotonía y cadlág):} Justificación analítica rigurosa del crecimiento no decreciente basada en la inclusión de eventos medibles $\set{X \le a} \subseteq \set{X \le b}$ y convergencia por arriba de límites por la derecha.
		\end{itemize}
	\end{block}
\end{frame}

% SLIDE 16: Problemas --- Nivel Operativo
\begin{frame}{Problemas del Libro Maestro --- Nivel Operativo}
	\begin{block}{Aplicaciones en Ingeniería y Ciberseguridad}
		\small
		\begin{itemize}
			\itemsep=0.08cm
			\item \textbf{Problema 3.2.4 (Control de calidad en obleas de silicio):} Modelado a trozos del número de soldaduras defectuosas por oblea ($X \in \set{0,1,2,3}$) y evaluación directa de la probabilidad operativa de rechazo industrial: $P(X \ge 2) = 1 - F_X(1) = 0.30$.
			\item \textbf{Problema 3.2.5 (Pruebas de penetración e intentos geométricos):} Deducción en forma cerrada de la CDF de ensayos hasta el primer éxito $F_N(n) = 1 - (1-p)^n$ y determinación de la cuota mínima de intentos de auditoría ($n^* = 14$ para $p=0.20$ con $95\%$ de certeza).
			\item \textbf{Problema 3.2.6 (Aprovisionamiento en centros de datos):} Determinación por el operador ínfimo de la mediana diaria de demanda ($q_{0.50} = 30$ servidores) y el cuantil de seguridad de orden superior ($q_{0.85} = 40$ servidores para SLA del $85\%$).
		\end{itemize}
	\end{block}
\end{frame}

% SLIDE 17: Problemas --- Nivel Analítico
\begin{frame}{Problemas del Libro Maestro --- Nivel Analítico}
	\begin{block}{Demostraciones Teóricas y Transformaciones de Variables}
		\small
		\begin{itemize}
			\itemsep=0.1cm
			\item \textbf{Problema 3.2.7 (Diferencias sobre enteros):} Demostración formal de que para variables con soporte en $\mathbb{Z}$, el intervalo cerrado equivale exactamente a la diferencia retardada:
			\begin{align*}
				P(a \le X \le b) = F_X(b) - F_X(a - 1) \quad \text{y} \quad f_X(k) = \Delta F_X(k).
			\end{align*}
			\item \textbf{Problema 3.2.8 (Transformaciones monótonas):} Demostración de invarianza estructural de escalones al transformar variables $Y = g(X)$ creciente ($F_Y(y) = F_X(g^{-1}(y))$) y deducción formal para funciones decrecientes $Z = h(X)$:
			\begin{align*}
				F_Z(z) = 1 - \lim_{u \to (h^{-1}(z))^-} F_X(u).
			\end{align*}
		\end{itemize}
	\end{block}
\end{frame}

% SLIDE 18: Problemas --- Nivel Desafiante
\begin{frame}{Problemas del Libro Maestro --- Nivel Desafiante}
	\begin{block}{Teoría de Distribuciones y Cotas Bivariadas Extremas}
		\small
		\begin{itemize}
			\itemsep=0.1cm
			\item \textbf{Problema 3.2.9 (Distribuciones generalizadas y escalones de Heaviside):} Representación de la CDF discreta pura como combinación de escalones $F(x) = \sum p_k H(x - x_k)$ y demostración topológica en Sobolev-Schwartz de que su derivada distribucional reproduce impulsos de Dirac: $f(x) = \sum p_k \delta(x - x_k)$.
			\item \textbf{Problema 3.2.10 (Cotas conjuntas de Fréchet-Hoeffding en enrutadores):} Deducción del cálculo de masa en rectángulos por segundas diferencias y demostración de las cotas universales que confinan toda CDF conjunta entre sus marginales:
			\begin{align*}
				\max\left( 0, F_X(x) + F_Y(y) - 1 \right) \le F_{X,Y}(x, y) \le \min\left( F_X(x), F_Y(y) \right).
			\end{align*}
		\end{itemize}
	\end{block}
\end{frame}

% SLIDE 19: Conclusiones y Síntesis
\begin{frame}{Conclusiones --- La CDF como Pilar del Análisis Discreto}
	\begin{block}{Resumen de Competencias Adquiridas}
		\small
		\begin{enumerate}
			\itemsep=0.06cm
			\item \textbf{Equivalencia Estructural:} La Función de Distribución Acumulada $F(x)$ encapsula la totalidad de la información de un experimento aleatorio con igual rigor y unicidad que la PMF $f(x)$.
			\item \textbf{Geometría de Escalones:} En el caso discreto, $F(x)$ se caracteriza por tramos constantes (mesetas) intercalados por saltos verticales exactos en los puntos del soporte $\mathcal{S}_X$.
			\item \textbf{Inversión y Operaciones:} El operador de diferencia hacia atrás $f(x) = F(x) - F(x^-)$ y las fórmulas de intervalos permiten evaluar riesgos sin recurrir a sumatorias tediosas.
		\end{enumerate}
	\end{block}
	\vspace{-0.15cm}
	\begin{alertblock}{Verificación Computacional}
		\footnotesize
		La simulación Monte Carlo confirmó que la distribución empírica (ECDF) converge uniformemente a la CDF teórica, validando el Teorema Fundamental de Glivenko-Cantelli.
	\end{alertblock}
\end{frame}

% SLIDE 20: Puente Didáctico
\begin{frame}{Puente Didáctico --- Hacia la Esperanza Matemática y Momentos}
	\begin{columns}[T]
		\begin{column}{0.48\textwidth}
			\begin{block}{Lo que hemos logrado (03.01 y 03.02)}
				\small
				\begin{itemize}
					\itemsep=0.06cm
					\item Dominamos la descripción \emph{completa} y micro-analítica de variables discretas mediante $f(x)$ (masas puntuales) y $F(x)$ (acumulaciones escalonadas).
					\item Podemos responder consultas exactas, de intervalos, cuantiles y colas superiores/inferiores.
				\end{itemize}
			\end{block}
		\end{column}
		\begin{column}{0.48\textwidth}
			\begin{block}{Lo que sigue: Sección 03.03}
				\small
				\begin{itemize}
					\itemsep=0.06cm
					\item ¿Cómo resumir toda la distribución en un solo número representativo de tendencia central y dispersión?
					\item Introduciremos el operador de \textbf{Esperanza Matemática} $\mathbb{E}[X] = \sum x f(x)$, la \textbf{Varianza} $\text{Var}(X)$ y la teoría general de momentos y funciones generadoras.
				\end{itemize}
			\end{block}
		\end{column}
	\end{columns}
	\vspace{0.2cm}
	\begin{center}
		\textbf{\textcolor{TecAzul}{\Large ¡Felicidades por completar la Sección 03.02!}}
	\end{center}
\end{frame}

\end{document}
